\documentclass[aspectratio=169]{beamer}

% Font encoding and input
\usepackage[T1]{fontenc}
\usepackage[utf8]{inputenc}
\usepackage{microtype}

% Core packages
\usepackage{graphicx}
\usepackage{url}
\usepackage{listings}
\usepackage{booktabs}
\usepackage{tabularx}
\usepackage{xcolor}
\usepackage{amsmath}
\usepackage{amsfonts}
\usepackage{amssymb}

\usepackage[spanish]{babel} %use your language

\setbeamertemplate{bibliography item}[text]

\usetheme{UTN-BHI}

% El tema UTN-BHI está diseñado para 4:3. Como esta presentación es 16:9,
% se centra el fondo a la altura de la página en lugar de estirarlo, para que
% el logo no quede deformado.
\usebackgroundtemplate{%
  \hbox to \paperwidth{\hss\includegraphics[height=\paperheight]{theme/utn-bhi-page}\hss}}

% Espacio reservado al pie para que el texto no se superponga con el logo.
\setbeamertemplate{footline}{\rule{0pt}{0.9cm}}

% Tablas un poco más compactas, para que las más largas entren en la diapositiva.
\usepackage{etoolbox}
\renewcommand{\arraystretch}{0.95}
\AtBeginEnvironment{tabular}{\footnotesize}
% Sin justificado en las celdas de tipo p/X: evita los huecos entre palabras y
% los cortes de palabra raros en las columnas angostas.
\usepackage{ragged2e}
\makeatletter
\g@addto@macro\@arrayparboxrestore{\RaggedRight\hyphenpenalty=5000}
\makeatother

% Add a section slide.
% Remove it if you prefer that option.
% \AtBeginSection[]{
%   \begin{frame}
%   \vfill
%   \centering
%   \begin{beamercolorbox}[sep=8pt,center]{title}
%     \usebeamerfont{title}\insertsectionhead\par%
%   \end{beamercolorbox}
%   \vfill
%   \end{frame}
% }

% Enhanced code style
\lstdefinestyle{myC++Style}{
  language=C++,
  basicstyle=\ttfamily\small,
  keywordstyle=\color{blue}\bfseries,
  commentstyle=\color{green!60!black},
  stringstyle=\color{red},
  numberstyle=\tiny\color{gray},
  showspaces=false,
  showstringspaces=false,
  breaklines=true,
  numbers=left,
  frame=single,
  backgroundcolor=\color{gray!10},
  captionpos=b
}

% Estilo por defecto para todos los listados: el de arriba, pero en C y en un
% cuerpo que entra en las columnas de las diapositivas comparativas.
\lstset{
  style=myC++Style,
  language=C,
  basicstyle=\ttfamily\scriptsize,
  columns=flexible,
  numbersep=4pt,
  xleftmargin=1.2em,
  aboveskip=0.6em,
  belowskip=0.3em
}

\urldef{\email}\path|jiparraguirre@frbb.utn.edu.ar|
\title{Dos formas de escribir el mismo programa en C}
\author[J. Iparraguirre] {Javier Iparraguirre}

\institute{
Universidad Tecnológica Nacional\\
11 de abril 461, Bahía Blanca, Argentina\\
\email\\
\url{http://www.frbb.utn.edu.ar/}}



\begin{document}

% ------------------------------------------------------------
\begin{frame}
    \titlepage
\end{frame}

% ------------------------------------------------------------
\begin{frame}{Objetivo de la clase}
    \begin{itemize}
        \item Partimos de un \textbf{mismo enunciado} y una \textbf{misma solución funcional}.
        \item Vamos a comparar dos formas de escribirla:
        \begin{itemize}
            \item[\textbf{(1)}] Estilo \textbf{clásico}: el que normalmente alcanza para aprobar un examen.
            \item[\textbf{(2)}] Estilo \textbf{moderno}: el que se espera en código de uso real.
        \end{itemize}
        \item La idea \textbf{no} es que un estilo sea ``correcto'' y el otro ``incorrecto''.
        \item La idea es entender \textbf{qué problema resuelve cada decisión de diseño}.
    \end{itemize}
\end{frame}

% ------------------------------------------------------------
\begin{frame}{El problema que vamos a resolver}
    \small
    Escribir un programa que:
    \begin{enumerate}
        \item Genere un archivo de texto con $N$ enteros aleatorios en $[1000, 2000]$.
        \item Cargue esos $N$ números desde el archivo a un arreglo.
        \item Ordene el arreglo por el método de \textbf{inserción}.
        \item Imprima el arreglo ordenado.
        \item Todo integrado desde \texttt{main()}, pidiendo $N$ al usuario.
    \end{enumerate}

    \vspace{0.5em}
    \begin{block}{Por qué este ejercicio es un buen caso de estudio}
        Combina \textbf{archivos}, \textbf{arreglos}, \textbf{memoria dinámica} y
        \textbf{entrada de usuario}: los cuatro lugares donde más se nota la diferencia
        entre "que compile" y "que sea robusto".
    \end{block}
\end{frame}

% ------------------------------------------------------------
\begin{frame}{Dos estilos, mismo comportamiento}
    \begin{columns}[T]
        \column{0.5\textwidth}
        \begin{block}{Estilo clásico}
            \begin{itemize}
                \item \texttt{int} para todo (incluso tamaños)
                \item \texttt{scanf} directo
                \item Poca verificación de errores
                \item Literales ``mágicos''
            \end{itemize}
        \end{block}

        \column{0.5\textwidth}
        \begin{block}{Estilo moderno}
            \begin{itemize}
                \item \texttt{size\_t}, \texttt{bool}
                \item \texttt{fgets} + \texttt{strtol}
                \item Toda función que puede fallar, se verifica
                \item Constantes con nombre
            \end{itemize}
        \end{block}
    \end{columns}

    \vspace{1em}
    \begin{center}
        \textit{Ambos programas producen exactamente el mismo resultado.}
    \end{center}
\end{frame}

% ------------------------------------------------------------
\begin{frame}{Tabla comparativa general}
    \small
    \begin{table}
        \begin{tabular}{p{3.3cm}p{4.2cm}p{4.2cm}}
            \toprule
            \textbf{Aspecto} & \textbf{Estilo clásico} & \textbf{Estilo moderno} \\
            \midrule
            Tipos para tamaños & \texttt{int} & \texttt{size\_t} \\
            Valores booleanos & \texttt{int} (0/1) & \texttt{bool} (\texttt{stdbool.h}) \\
            Lectura de enteros & \texttt{scanf("\%d", \ldots)} & \texttt{fgets} + \texttt{strtol} \\
            Verificación de I/O & Parcial & En cada llamada \\
            Visibilidad de funciones & Global (implícita) & \texttt{static} cuando aplica \\
            Constantes & Literales en el código & \texttt{\#define} con nombre \\
            Reserva de memoria & \texttt{sizeof(int)} & \texttt{sizeof *puntero} \\
            \bottomrule
        \end{tabular}
    \end{table}
\end{frame}

% ------------------------------------------------------------
\begin{frame}[fragile]{1. Tipos para representar tamaños}
    \begin{columns}[T]
        \column{0.48\textwidth}
        \textbf{Clásico}
        \begin{lstlisting}
void ordenarInsercion(
    int arreglo[], int n)
{
    for (int i = 1; i < n; i++) {
        ...
    }
}
        \end{lstlisting}

        \column{0.48\textwidth}
        \textbf{Moderno}
        \begin{lstlisting}
static void ordenarInsercion(
    int *arreglo, size_t n)
{
    for (size_t i = 1; i < n; i++) {
        ...
    }
}
        \end{lstlisting}
    \end{columns}

    \vspace{0.5em}
    \begin{alertblock}{¿Por qué importa?}
        Un tamaño o una cantidad de elementos \textbf{nunca es negativo}.
        \texttt{size\_t} es un tipo \textbf{sin signo} que expresa esa garantía
        y evita comparaciones ambiguas entre \texttt{int} con signo y valores sin signo.
    \end{alertblock}
\end{frame}

% ------------------------------------------------------------
\begin{frame}[fragile]{2. Comunicar éxito o fracaso}
    \begin{columns}[T]
        \column{0.48\textwidth}
        \textbf{Clásico}
        \begin{lstlisting}
void generarArchivo(int n)
{
    FILE *f = fopen(FILENAME,"w");
    if (f == NULL) {
        printf("Error...\n");
        exit(1);
    }
    ...
}
        \end{lstlisting}

        \column{0.48\textwidth}
        \textbf{Moderno}
        \begin{lstlisting}
static bool generarArchivo(size_t n)
{
    FILE *f = fopen(FILENAME,"w");
    if (!f) {
        fprintf(stderr, "Error...\n");
        return false;
    }
    ...
    return true;
}
        \end{lstlisting}
    \end{columns}

    \vspace{0.5em}
    \begin{alertblock}{¿Por qué importa?}
        \texttt{exit(1)} corta el programa \textbf{desde adentro} de una función auxiliar,
        sin darle a quien la llamó la oportunidad de decidir qué hacer.
        Devolver \texttt{bool} deja esa decisión en manos de \texttt{main}.
    \end{alertblock}
\end{frame}

% ------------------------------------------------------------
\begin{frame}[fragile]{3. Verificación de cada operación de E/S}
    \begin{columns}[T]
        \column{0.48\textwidth}
        \textbf{Clásico}
        \begin{lstlisting}
for (int i = 0; i < n; i++) {
    fscanf(f, "%d", &arreglo[i]);
}
        \end{lstlisting}
        \vspace{1em}
        Si el archivo tiene menos números de los esperados, el arreglo queda
        con \textbf{basura} sin que nadie se entere.

        \column{0.48\textwidth}
        \textbf{Moderno}
        \begin{lstlisting}
for (size_t i = 0; i < n; i++) {
    if (fscanf(f, "%d",
               &arreglo[i]) != 1) {
        fprintf(stderr, "Error...\n");
        fclose(f);
        return false;
    }
}
        \end{lstlisting}
    \end{columns}

    \vspace{0.3em}
    \begin{alertblock}{Idea clave}
        \texttt{fscanf} devuelve la cantidad de campos que logró leer.
        Ignorar ese valor de retorno es ignorar información que la función
        ya nos está dando.
    \end{alertblock}
\end{frame}

% ------------------------------------------------------------
\begin{frame}[fragile]{4. Leer un número desde teclado}
    \begin{columns}[T]
        \column{0.48\textwidth}
        \textbf{Clásico}
        \begin{lstlisting}
int n;
printf("Ingrese N: ");
scanf("%d", &n);
        \end{lstlisting}
        \vspace{1em}
        \small Si el usuario escribe \texttt{"abc"}, el comportamiento
        es indefinido: \texttt{n} puede quedar sin inicializar.

        \column{0.48\textwidth}
        \textbf{Moderno}
        \begin{lstlisting}
char linea[64];
fgets(linea, sizeof linea, stdin);

char *endptr;
long valor = strtol(linea,
                     &endptr, 10);

if (endptr == linea || valor <= 0)
    /* entrada invalida */
        \end{lstlisting}
    \end{columns}

    \vspace{0.3em}
    \begin{alertblock}{Idea clave}
        \texttt{fgets} siempre lee una línea completa (sin desbordar el buffer);
        \texttt{strtol} permite \textbf{detectar} si lo que se escribió era
        realmente un número.
    \end{alertblock}
\end{frame}

% ------------------------------------------------------------
\begin{frame}[fragile]{5. \texttt{const} correctness}
    \begin{columns}[T]
        \column{0.48\textwidth}
        \textbf{Clásico}
        \begin{lstlisting}
void imprimirArreglo(
    int arreglo[], int n)
        \end{lstlisting}

        \column{0.48\textwidth}
        \textbf{Moderno}
        \begin{lstlisting}
static void imprimirArreglo(
    const int *arreglo, size_t n)
        \end{lstlisting}
    \end{columns}

    \vspace{1em}
    \begin{alertblock}{¿Por qué importa?}
        La firma con \texttt{const} es una \textbf{promesa documentada en el
        propio código}: esta función \textit{no puede} modificar el arreglo.
        El compilador la hace cumplir; no depende de leer el cuerpo de la función
        para confiar en eso.
    \end{alertblock}
\end{frame}

% ------------------------------------------------------------
\begin{frame}[fragile]{6. Reserva de memoria dinámica}
    \begin{columns}[T]
        \column{0.48\textwidth}
        \textbf{Clásico}
        \begin{lstlisting}
int *arreglo =
    (int *) malloc(n * sizeof(int));
        \end{lstlisting}

        \column{0.48\textwidth}
        \textbf{Moderno}
        \begin{lstlisting}
int *arreglo =
    malloc(n * sizeof *arreglo);
        \end{lstlisting}
    \end{columns}

    \vspace{1em}
    \begin{alertblock}{¿Por qué importa?}
        Si el día de mañana \texttt{arreglo} cambia de \texttt{int *} a
        \texttt{double *}, la versión con \texttt{sizeof *arreglo} sigue
        siendo correcta sin tocar esa línea. La versión con
        \texttt{sizeof(int)} queda desincronizada del tipo real.
    \end{alertblock}
\end{frame}

% ------------------------------------------------------------
\begin{frame}[fragile]{7. Constantes con nombre}
    \begin{columns}[T]
        \column{0.48\textwidth}
        \textbf{Clásico}
        \begin{lstlisting}
int numero = 1000 + rand() % 1001;
        \end{lstlisting}

        \column{0.48\textwidth}
        \textbf{Moderno}
        \begin{lstlisting}
#define RANGO_MIN 1000
#define RANGO_MAX 2000
...
int numero = RANGO_MIN +
    rand() % (RANGO_MAX - RANGO_MIN + 1);
        \end{lstlisting}
    \end{columns}

    \vspace{1em}
    \begin{alertblock}{¿Por qué importa?}
        El \texttt{1001} del estilo clásico es un ``número mágico'': funciona,
        pero no dice por qué es \texttt{1001} y no otro valor. Con constantes
        nombradas, el rango se puede cambiar en \textbf{un solo lugar}.
    \end{alertblock}
\end{frame}

% ------------------------------------------------------------
\begin{frame}[fragile]{8. Visibilidad de las funciones: \texttt{static}}
    \begin{columns}[T]
        \column{0.48\textwidth}
        \textbf{Clásico}
        \begin{lstlisting}
void ordenarInsercion(
    int arreglo[], int n)
        \end{lstlisting}
        \small Visible desde cualquier otro archivo \texttt{.c} del proyecto.

        \column{0.48\textwidth}
        \textbf{Moderno}
        \begin{lstlisting}
static void ordenarInsercion(
    int *arreglo, size_t n)
        \end{lstlisting}
        \small Solo visible dentro de este archivo.
    \end{columns}

    \vspace{1em}
    \begin{alertblock}{¿Por qué importa?}
        Cuando un programa crece a varios archivos, marcar como \texttt{static}
        las funciones que son ``detalles internos'' evita colisiones de nombres
        y deja claro cuál es la interfaz pública del módulo.
    \end{alertblock}
\end{frame}

% ------------------------------------------------------------
\begin{frame}{Un paso más: C11 y C23}
    \begin{itemize}
        \item Hasta ahora comparamos \textbf{C99 clásico} vs. \textbf{C99 moderno}:
              dos formas de escribir lo mismo dentro del \textbf{mismo estándar}.
        \item El lenguaje C también evolucionó: \textbf{C11} (2011) y
              \textbf{C23} (2024) agregaron herramientas nuevas al lenguaje.
        \item Vamos a ver cuáles de esas herramientas se aplican directamente
              a nuestro ejercicio, y por qué.
    \end{itemize}

    \vspace{1em}
    \begin{alertblock}{Antes de avanzar: soporte de compiladores}
        C23 es muy reciente. GCC lo soporta desde la versión 13 (con la bandera
        \texttt{-std=c2x} o, en versiones más nuevas, \texttt{-std=c23}) y Clang
        desde la versión 16--18 según la característica. \textbf{Siempre
        verifiquen la versión del compilador} antes de usar estas
        características en un trabajo práctico.
    \end{alertblock}
\end{frame}

% ------------------------------------------------------------
\begin{frame}{Historial cronológico de versiones de C (1/2)}
    \small
    \begin{table}
        \begin{tabularx}{\textwidth}{l p{2.8cm} c X}
            \toprule
            \textbf{Versión} & \textbf{Nombre del estándar} & \textbf{Año} & \textbf{Hito clave / características añadidas} \\
            \midrule
            \textbf{K\&R C} & Pre-estándar & 1978 & Escrito por Brian Kernighan y Dennis Ritchie; base original e informal del lenguaje. \\
            \addlinespace
            \textbf{C89} & ANSI X3.159-1989 & 1989 & Primera versión oficial estandarizada, para mejorar la portabilidad entre plataformas. \\
            \addlinespace
            \textbf{C90} & ISO/IEC 9899:1990 & 1990 & Idéntica a C89, adoptada formalmente por ISO a nivel global. \\
            \addlinespace
            \textbf{C95} & ISO/IEC 9899/AMD1 & 1995 & Soporte para caracteres multibyte y anchos, para software internacional. \\
            \addlinespace
            \textbf{C99} & ISO/IEC 9899:1999 & 1999 & Funciones \texttt{inline}, arreglos de longitud variable, \texttt{long long int}, comentarios \texttt{//}. \\
            \bottomrule
        \end{tabularx}
    \end{table}
\end{frame}

% ------------------------------------------------------------
\begin{frame}{Historial cronológico de versiones de C (2/2)}
    \small
    \begin{table}
        \begin{tabularx}{\textwidth}{l p{2.8cm} c X}
            \toprule
            \textbf{Versión} & \textbf{Nombre del estándar} & \textbf{Año} & \textbf{Hito clave / características añadidas} \\
            \midrule
            \textbf{C11} & ISO/IEC 9899:2011 & 2011 & Hilos (multithreading), operaciones atómicas y Unicode. \\
            \addlinespace
            \textbf{C17} & ISO/IEC 9899:2018 & 2018 & Mantenimiento y corrección de errores; sin características nuevas. \\
            \addlinespace
            \textbf{C23} & ISO/IEC 9899:2024 & 2024 & Estándar más reciente ratificado: \texttt{nullptr}, \texttt{bool}, \texttt{typeof}, atributos. \\
            \addlinespace
            \textbf{C2Y} & En desarrollo & Próximo & Borrador experimental actual para la próxima especificación oficial. \\
            \bottomrule
        \end{tabularx}
    \end{table}

    \vspace{1em}
    \small Las próximas diapositivas profundizan en \textbf{C11} y \textbf{C23},
    las dos versiones más relevantes para nuestro ejercicio.
\end{frame}

% ------------------------------------------------------------
\begin{frame}{Novedades de C11 (2011)}
    \small
    \begin{table}
        \begin{tabular}{p{4.3cm}p{7.5cm}}
            \toprule
            \textbf{Característica} & \textbf{Para qué sirve} \\
            \midrule
            \texttt{\_Static\_assert} & Verificar condiciones \textbf{en tiempo de compilación} \\
            \texttt{\_Generic} & Elegir código según el tipo de un argumento (``sobrecarga'' simple) \\
            \texttt{\_Noreturn} & Documentar que una función nunca retorna (ej. un wrapper de \texttt{exit}) \\
            Uniones/estructuras anónimas & Acceder a miembros sin nombrar la unión/estructura contenedora \\
            \texttt{\_Alignas} / \texttt{\_Alignof} & Controlar y consultar el alineamiento en memoria de un tipo \\
            \texttt{<threads.h>}, \texttt{<stdatomic.h>} & Hilos y operaciones atómicas estandarizadas (antes dependían del SO) \\
            \bottomrule
        \end{tabular}
    \end{table}

    \vspace{0.3em}
    \small De todas estas, la que más directamente mejora nuestro ejercicio
    es \texttt{\_Static\_assert}.
\end{frame}

% ------------------------------------------------------------
\begin{frame}[fragile]{C11 aplicado: \texttt{static\_assert}}
    \begin{lstlisting}
static_assert(RANGO_MAX > RANGO_MIN,
    "RANGO_MAX debe ser mayor que RANGO_MIN");
    \end{lstlisting}

    \vspace{0.5em}
    \begin{alertblock}{¿Por qué importa?}
        Esta verificación se evalúa \textbf{cuando se compila} el programa,
        no cuando se ejecuta. Si alguien invierte los valores de
        \texttt{RANGO\_MIN} y \texttt{RANGO\_MAX} por error, el programa
        \textbf{ni siquiera compila} — no hace falta ejecutarlo ni escribir
        un test para descubrir el error.
    \end{alertblock}

    \vspace{0.3em}
    \small \texttt{<assert.h>} expone la macro \texttt{static\_assert} desde
    C11 (por debajo, expande a \texttt{\_Static\_assert}). En C23 pasó a ser
    directamente una palabra clave del lenguaje.
\end{frame}

% ------------------------------------------------------------
\begin{frame}{Novedades de C23 (2024)}
    \small
    \begin{table}
        \begin{tabular}{p{4.3cm}p{7.5cm}}
            \toprule
            \textbf{Característica} & \textbf{Para qué sirve} \\
            \midrule
            \texttt{bool}/\texttt{true}/\texttt{false} & Palabras clave del lenguaje: ya no se necesita \texttt{<stdbool.h>} \\
            \texttt{nullptr} & Puntero nulo con tipo propio, en vez de reusar \texttt{NULL}/\texttt{0} \\
            \texttt{constexpr} & Constantes evaluadas en tiempo de compilación, con tipo (mejora a \texttt{\#define}) \\
            \texttt{typeof} & Obtener el tipo de una expresión para usarlo en otra declaración \\
            \texttt{auto} & Inferencia de tipo a partir del inicializador \\
            \texttt{[[atributos]]} & Anotaciones estándar como \texttt{[[nodiscard]]}, \texttt{[[deprecated]]} \\
            \texttt{\#embed} & Incluir un archivo binario directamente como datos en el código \\
            \bottomrule
        \end{tabular}
    \end{table}

    \vspace{0.3em}
    \small De estas, \texttt{bool}/\texttt{true}/\texttt{false},
    \texttt{nullptr}, \texttt{constexpr} y \texttt{[[nodiscard]]} se aplican
    directamente a nuestro ejercicio.
\end{frame}

% ------------------------------------------------------------
\begin{frame}[fragile]{C23 aplicado: \texttt{bool} y \texttt{nullptr}}
    \begin{columns}[T]
        \column{0.48\textwidth}
        \textbf{Moderno (C99)}
        \begin{lstlisting}
#include <stdbool.h>
...
FILE *f = fopen(FILENAME, "w");
if (f == NULL) {
    ...
}
        \end{lstlisting}

        \column{0.48\textwidth}
        \textbf{C23}
        \begin{lstlisting}
/* sin stdbool.h: bool es
   palabra clave */
FILE *f = fopen(FILENAME, "w");
if (f == nullptr) {
    ...
}
        \end{lstlisting}
    \end{columns}

    \vspace{0.5em}
    \begin{alertblock}{¿Por qué importa?}
        \texttt{NULL} en C99 suele ser simplemente \texttt{((void*)0)} o
        incluso \texttt{0}: se puede confundir con un entero. \texttt{nullptr}
        tiene su propio tipo (\texttt{nullptr\_t}) y \textbf{solo} puede
        convertirse a puntero o a \texttt{bool} — el compilador detecta más
        errores de tipos.
    \end{alertblock}
\end{frame}

% ------------------------------------------------------------
\begin{frame}[fragile]{C23 aplicado: \texttt{constexpr} en vez de \texttt{\#define}}
    \begin{columns}[T]
        \column{0.48\textwidth}
        \textbf{Moderno (C99)}
        \begin{lstlisting}
#define RANGO_MIN 1000
#define RANGO_MAX 2000
        \end{lstlisting}

        \column{0.48\textwidth}
        \textbf{C23}
        \begin{lstlisting}
constexpr int RANGO_MIN = 1000;
constexpr int RANGO_MAX = 2000;
        \end{lstlisting}
    \end{columns}

    \vspace{0.5em}
    \begin{alertblock}{¿Por qué importa?}
        Una macro de preprocesador \textbf{no tiene tipo} y es solo un
        reemplazo de texto antes de compilar: el depurador no la ve, y el
        compilador no puede avisar si se usa con el tipo equivocado.
        \texttt{constexpr} es una variable real, con tipo, cuyo valor el
        compilador garantiza que se calculó en tiempo de compilación.
    \end{alertblock}
\end{frame}

% ------------------------------------------------------------
\begin{frame}[fragile]{C23 aplicado: \texttt{[[nodiscard]]}}
    \begin{lstlisting}
[[nodiscard]] static bool generarArchivo(size_t n);
    \end{lstlisting}

    \vspace{0.3em}
    Si alguien llama a la función e ignora el resultado:
    \begin{lstlisting}
generarArchivo(n);   /* se ignora el valor de retorno */
    \end{lstlisting}

    \vspace{0.3em}
    el compilador avisa \textbf{en tiempo de compilación}:
    \begin{lstlisting}[language=bash]
warning: ignoring return value of 'generarArchivo',
declared with attribute 'nodiscard' [-Wunused-result]
    \end{lstlisting}

    \begin{alertblock}{¿Por qué importa?}
        En el estilo moderno (C99) ignorar un \texttt{bool} de retorno es un
        error \textbf{silencioso}: compila sin quejas. \texttt{[[nodiscard]]}
        convierte ese descuido en un warning visible, verificado por el
        compilador y no por la disciplina de quien programa.
    \end{alertblock}
\end{frame}

% ------------------------------------------------------------
\begin{frame}[fragile]{Código completo: dónde encontrarlo}
    \begin{itemize}
        \item \texttt{codigo/clasico.c} \quad — programa completo, estilo clásico (C99).
        \item \texttt{codigo/moderno.c} \quad — programa completo, estilo moderno (C99).
        \item \texttt{codigo/c23.c} \quad — la versión moderna, incorporando
              \texttt{static\_assert} (C11) y \texttt{bool}/\texttt{nullptr}/
              \texttt{constexpr}/\texttt{[[nodiscard]]} (C23).
    \end{itemize}

    \vspace{1em}
    \begin{block}{Para la clase práctica}
        \begin{lstlisting}[language=bash]
gcc -std=c99 -Wall -Wextra -o clasico  clasico.c
gcc -std=c99 -Wall -Wextra -o moderno  moderno.c
gcc -std=c2x -Wall -Wextra -o c23      c23.c
# En GCC 15+ y Clang 18+ tambien funciona -std=c23
        \end{lstlisting}
        Comparen \textbf{cuántos warnings} tira \texttt{-Wall -Wextra} en
        cada caso, y qué errores detecta \texttt{c23.c} \textbf{antes} de
        ejecutar el programa (\texttt{static\_assert}) que en las otras dos
        versiones solo se notarían en tiempo de ejecución o ni se notarían.
    \end{block}
\end{frame}

% ------------------------------------------------------------
\begin{frame}{¿Cuándo usar cada estilo?}
    \begin{columns}[T]
        \column{0.48\textwidth}
        \begin{block}{Estilo clásico}
            \begin{itemize}
                \item Ejercicios y exámenes con tiempo acotado
                \item Prototipos rápidos de un solo uso
                \item Cuando el enunciado no exige robustez
            \end{itemize}
        \end{block}

        \column{0.48\textwidth}
        \begin{block}{Estilo moderno}
            \begin{itemize}
                \item Cualquier código que otra persona vaya a mantener
                \item Software que va a producción
                \item Trabajos prácticos evaluados por calidad, no solo por ``funciona''
            \end{itemize}
        \end{block}
    \end{columns}

    \vspace{1em}
    \begin{center}
        \textbf{No es una cuestión de gusto: es elegir la herramienta según el contexto.}
    \end{center}
\end{frame}

% ------------------------------------------------------------
\begin{frame}{Actividad para la clase}
    \begin{enumerate}
        \item Tomen \texttt{codigo/clasico.c} y compílenlo con \texttt{-Wall -Wextra}.
        \item Identifiquen al menos \textbf{tres} advertencias del compilador.
        \item Para cada una, expliquen qué principio del estilo moderno la evita.
        \item Reescriban la función \texttt{cargarArchivo} aplicando esos principios,
              \textbf{sin mirar} \texttt{moderno.c}.
        \item Comparen su solución con la versión moderna del repositorio.
        \item \textbf{Extra:} en \texttt{codigo/c23.c}, cambien a propósito
              \texttt{RANGO\_MIN} y \texttt{RANGO\_MAX} para que el rango sea
              inválido. Compilen de nuevo y observen qué mensaje da
              \texttt{static\_assert} \textbf{antes} de que el programa corra.
    \end{enumerate}
\end{frame}

% ------------------------------------------------------------
\begin{frame}{Resumen}
    \begin{itemize}
        \item El estilo moderno no agrega funcionalidad: agrega \textbf{garantías}.
        \item Cada cambio (tipos, verificación de errores, \texttt{const}, nombres)
              responde a una pregunta concreta: \textit{``¿qué puede salir mal y cómo lo detecto?''}
        \item Escribir C ``a la clásica'' está bien para aprender los conceptos
              base del lenguaje.
        \item Escribir C ``moderno'' es la base para escribir software confiable,
              algo central en Ingeniería Eléctrica cuando el código controla
              hardware real.
        \item C11 y C23 no cambian estas ideas: les dan al compilador
              \textbf{más información} (\texttt{static\_assert},
              \texttt{constexpr}, \texttt{[[nodiscard]]}) para que las
              garantías del estilo moderno se verifiquen automáticamente en
              vez de depender solo de la disciplina de quien programa.
    \end{itemize}
\end{frame}

\end{document}
